% !PRINTED 359
\weberpage{359}
Since $n$ is coprime to $Q$, one can set this number $= nr$,
where $r$ is a rational number coprime to $Q$. Accordingly every number $\eta'$ has the property
\weq{1}{\eta' \equiv r \pmod{Q},}
where $r$ is a rational integer and the congruence is to be
understood in the sense $n\eta' \equiv nr$.

The congruences (1) can be multiplied and divided
if by $r^{-1}$ one understands the integer satisfying the condition $r \cdot r^{-1} \equiv 1 \pmod{Q}$.

Conversely, an integral or fractional number $\eta'$ which
satisfies a condition (1) belongs to the group $O'$. For if
$n\eta' = \omega \equiv nr$ is a number in $[Q]$, then $\eta' = \omega : n$ is a quotient
of two such numbers.

3. In $O'$ there is now again contained a group $O_0$, which consists
of all numbers $\eta_0$ that satisfy the condition
\weq{2}{\eta_0 \equiv 1 \pmod{Q}}
and the numbers of $O_0$ likewise yield again, through multiplication
and division, numbers of $O_0$.

We thus have three kinds of number groups, each of which is contained
in the preceding one:

1. $O$ \gesperrt{contains all numbers in $\Omega$ prime to $Q$.}

2. $O'$ \gesperrt{contains the integral and fractional numbers of the order $[Q]$ prime to $Q$.}

3. $O_0$ \gesperrt{contains the number classes in $O'$ that are congruent to unity with respect to the modulus $Q$.}

The groups $O$, $O'$, $O_0$ are infinite \gesperrt{Abelian} groups.
But if one takes a complete system of residues of rational numbers coprime
 to $Q$, $r_1, r_2, \ldots, r_\mu$, then every number in $O'$ can be represented as the product
of one of these $r_i$ with a number in $O_0$, and one
can therefore set, in the mode of expression used in Bd.\,II, § 2:
\weq{3}{O' = r_1 O_0 + r_2 O_0 + \cdots + r_\mu O_0,}
and $\mu$ is the \gesperrt{index} of the divisor $O_0$ of $O$: % [sic] German source prints unprimed O.
\weq{4}{\mu = (O', O_0),}
hence a finite number, namely:
\weq{5}{\mu = \varphi(Q) = Q\, \Pi\left(1 - \frac{1}{q}\right),}
where $\varphi(Q)$ is the symbol customary in number theory for
the number of number classes with respect to the modulus $Q$ with numbers coprime to $Q$,
%CONT>>
